\documentclass[twoside, 12pt]{ctexart}

\usepackage{xcolor}
\usepackage{fontspec}
\usepackage{ruby} 

\usepackage{xcolor}
\usepackage{xparse}
\usepackage{calc}

\usepackage{relsize}
\usepackage{fancyhdr}

\usepackage{etoolbox,xstring}

\usepackage{geometry}
\usepackage[gb4e]{expex-glossonly}
\usepackage{expex}

\usepackage{booktabs}
\usepackage{array}

\usepackage{lipsum}

\usepackage{etoolbox}

\usepackage{booktabs,array}
\usepackage{caption}
% globally for all tables:
\captionsetup[table]{%
  name=,        % no “Table”
  labelsep=colon% “1: Caption” 
}

\geometry{a5paper, portrait, margin=2cm}

\pagestyle{fancy}
\fancypagestyle{firstpage}
{
    \fancyhead[L]{}    
    \fancyhead[R]{}
    \fancyfoot[L]{}
    \fancyfoot[C]{JGantts · 人カント}
}
\fancyhf{}
\fancyhead[LE]{画ひ言新}
\fancyhead[LO]{sitelen pi toki sin}
\fancyhead[RE]{Toki Sin Writing}
\fancyfoot[LO]{jan kanto}
\fancyfoot[RE]{人カント}
\fancyfoot[LE,RO]{\thepage}

\definecolor{primary}{RGB}{47, 128, 190}    % Muted Blue
\definecolor{secondary}{RGB}{192, 57, 43}   % Deep Red
\definecolor{accent}{RGB}{39, 174, 96}      % Cool Green
\definecolor{background}{RGB}{255, 255, 255} % Softer Light Gray
\definecolor{link}{RGB}{128, 0, 128}       % Regal Purple

\definecolor{pagedark}{RGB}{0, 0, 0}            %black
\definecolor{textdark}{RGB}{234, 234, 234}      %grey

\definecolor{pagelight}{RGB}{255, 255, 255}            %white
\definecolor{textlight}{RGB}{0, 0, 0}            %black

\ifnum\lightMode=1
  \colorlet{page}{pagelight}
  \colorlet{text}{textlight}
\else
  \colorlet{page}{pagedark}
  \colorlet{text}{textdark}
\fi

\pagecolor{page}
\color{text}

\setmainfont{NotoSerif}[
  Path = fonts/,
  Extension = .ttf,
  UprightFont = *-Light,
  ItalicFont = *-LightItalic,
  BoldFont = *-Bold,
  BoldItalicFont = *-BoldItalic,
]

\setsansfont{NotoSans}[
  Path = fonts/,
  Extension = .ttf,
  UprightFont = *-Light,
  ItalicFont = *-LightItalic,
  BoldFont = *-Bold,
  BoldItalicFont = *-BoldItalic,
]

\setCJKmainfont{NotoSerifJP}[
  Path = fonts/,
  Extension = .ttf,
  UprightFont = *-Regular,
  BoldFont = *-Black,
]

\newCJKfontfamily\cjksans{NotoSansJP}[
  Path = fonts/,
  Extension = .ttf,
  UprightFont = *-Regular,
  BoldFont = *-Black,
]

\NewDocumentCommand{\jgruby}{m O{}}{%
  \ruby{#1}{\upshape\sffamily{#2}}
}

\ExplSyntaxOn
\NewDocumentCommand\rby{>{\SplitList{~}}m}{\ProcessList{#1}{\rby_map}}
\NewDocumentCommand\rby_map{>{\SplitArgument{1}{=}}m}{{\rby_impl #1}}
\NewDocumentCommand\rby_impl{m m}{\IfNoValueTF{#2}{\jgruby{#1}}{\jgruby{#1}[#2]}}
\ExplSyntaxOff

\NewDocumentCommand{\srby}{m O{}}{%
  \mbox{\rby{#1}{#2}}
}

\def\maxvalue(#1,#2){\ifnum #1>#2 #1\else #2\fi}

\NewDocumentCommand{\englishundertokisin}{m m}{%
  \begin{minipage}[c]{\maxof{\widthof{#1}}{\widthof{#2}}}
    \centering{#1}\small\\\centering{#2}
  \end{minipage}%
}


\begin{document}
%\lipsum[1-20]



% {\Huge{}
% \noindent%
% \rby{人=jan カント=kanto}\\
% \\
% \begin{center}
% \rby{人=jan カント=kanto 因=tan}\\
% \rby{土=ma ひ=pi 山=nena アハラシヤ=apalasija}\\
% \end{center}
% \begin{center}
% \rby{人 カント=kanto 因}\\
% \rby{土 ひ 山=nena アハラシヤ=apalasija}\\
% \end{center}
% }

% \newpage
\sffamily{}
\thispagestyle{firstpage}
\vspace{1cm}

\begin{center}
    {\fontsize{48}{60}\bfseries\rmfamily\color{primary}\rby{画=sitelen ひ=pi 言=toki 新=sin}}\LARGE\\
\end{center}
\begin{center}
    {\Huge\rmfamily\textsf{}Toki Sin Writing}
\end{center}

\section[Overview]{\Huge{}Overview}
This document covers \rby{画=sitelen ひ=pi 言=toki 新=sin}, or toki sin writing.
There are three sets of characters used.\\
{{\textbf{\rby{画=sitelen ラシナ=lasina}} is a portion of the {\textbf{Latin Alphabet}}, the same letters used in English, French, German, etc.}\\
{{\textbf{\rby{画=sitelen 漢=kan}} is a portion of {\textbf{Chinese characters, 漢字 (hanzi/kanji)}}; with some particles being {\textbf{Japanese Hiragana}}.}\\
{{\textbf{\rby{画=sitelen 片=kata}} is simply {\textbf{Japanese Katakana}}.}\\

\begin{table}[ht]
  \centering
  \caption{The phrase 「\rby{言=toki 新=sin}」 in the three toki sin writing systems}
  \begin{tabular}{@{} m{6cm} l @{}}
    \toprule
    \textbf{Writing System} & \textbf{Representation} \\
    \midrule
    \rby{画=sitelen ラシナ=lasina} & {toki sin} \\
    \rby{画=sitelen 漢=kan}        & {言新}    \\
    \rby{画=sitelen 片=kata}       & {トキシン} \\
    \bottomrule
  \end{tabular}
\end{table}

\section[画ラシナ – Latin Characters]{
  \Huge\englishundertokisin
    {\rby{画=sitelen ラシナ=lasina}}
    {\small\normalfont{}Latin Characters}
}

\noindent
\begin{minipage}[t]{3.5cm}
\centering
\textbf{Consonants}\par\vspace{2pt}
\begin{tabular}{@{}ccc@{}}
\toprule
m & n &   \\
p & t & k \\
  & s &   \\
w & l & j \\
\bottomrule
\end{tabular}
\end{minipage}
\begin{minipage}[t]{2cm}
\centering
\textbf{Vowels}\par\vspace{2pt}
\begin{tabular}{@{}ccc@{}}
\toprule
i &   & u \\
e &   & o \\
  & a &   \\
\bottomrule
\end{tabular}
\end{minipage}\\
\vspace{0.5cm}\\

Characters called
\textbf{\rby{画=sitelen ラシナ=lasina}}, or \textbf{Latin characters},
are used to indicate pronunciation in \rby{画=sitelen ひ=pi 言=toki 新=sin} {\textit{(toki sin writing)}}.
They will typically be displayed above the corresponding 
\rby{画=sitelen 漢=kan} or \rby{画=sitelen 片=kata}.
They are mainly used in material intended for learners;
but they are also used for \rby{画=sitelen 漢=kan ひ=pi 多=mute 不=ala} {\textit{(uncommon \rby{画=sitelen 漢=kan})}}.

\newpage
\section[画漢 – Sitelen Kan]{
  \Huge\englishundertokisin
    {\rby{画=sitelen 漢=kan}}
    {\small\normalfont{}Kan Characters}
}

{
\begin{exe}
\ex[]{\label{foo}
\begingl
\gls \rby{人=jan り=li 食=moku え=e 果=kili 。}\\
\gla \sffamily{人} \sffamily{り} \sffamily{食} \sffamily{え} \sffamily{果}//
\glb jan li moku e kili//
\glc person PRED eat DO fruit//
\glft 'Person eat fruit.'//
\endgl
}
\end{exe}
}

In the above example, the \rby{画=sitelen 漢=kan} are used to write the words in toki sin.
The ruby text above the \rby{画=sitelen 漢=kan} indicates how to pronounce them.

\newpage
\section[画多 – Common Characters]{
  \Huge\englishundertokisin
    {\rby{画=sitelen 漢=kan 多=mute}}
    {\small\normalfont{}Common Kan Characters}
}

Here are the \rby{画=sitelen 漢=kan 多=mute},
common \rby{画=sitelen 漢=kan}.
\rby{画=sitelen} which aren't on this list will frequently have ruby text above them to assist in reading.\\

{
\Huge
\rby{私=mi 君=sina 其=ona}

\rby{因=tan 要=wile 来=kama 能=kin}

\rby{ら=la り=li ひ=pi}

\rby{え=e}

\rby{あ=a えん=en あぬ=anu 許=taso 内=insa}

\rby{不=ala 何=seme}

\rby{人=jan 于=lon 去=tawa 土=ma}
}

\newpage
\section[画片 – Sitelen Kata]{
  \Huge\englishundertokisin
    {\rby{画=sitelen 片=kata}}
    {\small\normalfont{}Kata Characters}
}

{
\Huge
\noindent
\rby{ア=a イ=i ウ=u エ=e オ=o}\\
  \rby{マ=ma ミ=mi ム=mu メ=me モ=mo}\\
  \rby{ナ=na ニ=ni ヌ=nu ネ=ne ノ=no}\\
  \rby{ハ=pa ヒ=pi フ=pu へ=pe ホ=po}\\
  \rby{タ=ta 　 ツ=tu テ=te ト=to}\\
  \rby{カ=ka キ=ki ク=ku ケ=ke コ=ko}\\
  \rby{サ=sa シ=si ス=su セ=se ソ=so}\\
  \rby{ラ=la リ=li ル=lu レ=le ロ=lo}\\
  \rby{ヤ=ya 　 ユ=yu イエ=ye ヨ=yo}\\
  \rby{ワ=wa ヰ=wi 　 ヱ=we 　}\\
  \rby{　 　 ン=n}\\
}

\newpage
\section[画漢 – Sitelen Kan]{
  \Huge\englishundertokisin
    {\rby{画=sitelen 漢=kan 全=ale}}
    {\small\normalfont{}Kan Characters, Exhaustive}
}
{\huge\noindent
\rby{あ=a}\hspace{0.3em}\rby{
獣=akesi}\hspace{0.3em}\rby{
不=ala}\hspace{0.3em}\rby{
集=alasa}\hspace{0.3em}\rby{
全=ale}\hspace{0.3em}\rby{
下=anpa}\hspace{0.3em}\rby{
化=ante}\hspace{0.3em}\rby{
あぬ=anu}\hspace{0.3em}\rby{
待=awen}\hspace{0.3em}\\
\rby{
え=e}\hspace{0.3em}\rby{
えん=en}\hspace{0.3em}\rby{
市=esun}\hspace{0.3em}\\
\rby{
物=ijo}\hspace{0.3em}\rby{
悪=ike}\hspace{0.3em}\rby{
具=ilo}\hspace{0.3em}\rby{
内=insa}\hspace{0.3em}\\
\rby{
汚=jaki}\hspace{0.3em}\rby{
人=jan}\hspace{0.3em}\rby{
黄=jelo}\hspace{0.3em}\rby{
有=jo}\hspace{0.3em}\\
\rby{
魚=kala}\hspace{0.3em}\rby{
音=kalama}\hspace{0.3em}\rby{
来=kama}\hspace{0.3em}\rby{
木=kasi}\hspace{0.3em}\rby{
能=ken}\hspace{0.3em}\rby{
使=kepeken}\hspace{0.3em}\rby{
果=kili}\hspace{0.3em}\rby{
又=kin}\hspace{0.3em}\rby{
切=kipisi}\hspace{0.3em}\rby{
石=kiwen}\hspace{0.3em}\rby{
粉=ko}\hspace{0.3em}\rby{
空=kon}\hspace{0.3em}\rby{
色=kule}\hspace{0.3em}\rby{
群=kulupu}\hspace{0.3em}\rby{
聞=kute}\hspace{0.3em}\\
\rby{
ら=la}\hspace{0.3em}\rby{
眠=lape}\hspace{0.3em}\rby{
青=laso}\hspace{0.3em}\rby{
首=lawa}\hspace{0.3em}\rby{
布=len}\hspace{0.3em}\rby{
冷=lete}\hspace{0.3em}\rby{
り=li}\hspace{0.3em}\rby{
小=lili}\hspace{0.3em}\rby{
糸=linja}\hspace{0.3em}\rby{
葉=lipu}\hspace{0.3em}\rby{
赤=loje}\hspace{0.3em}\rby{
ろ=lo}\hspace{0.3em}\rby{
于=lon}\hspace{0.3em}\rby{
手=luka}\hspace{0.3em}\rby{
見=lukin}\hspace{0.3em}\rby{
穴=lupa}\hspace{0.3em}\\
\rby{
土=ma}\hspace{0.3em}\rby{
母=mama}\hspace{0.3em}\rby{
貝=mani}\hspace{0.3em}\rby{
女=meli}\hspace{0.3em}\rby{
私=mi}\hspace{0.3em}\rby{
男=mije}\hspace{0.3em}\rby{
も=mo}\hspace{0.3em}\rby{
食=moku}\hspace{0.3em}\rby{
死=moli}\hspace{0.3em}\rby{
後=monsi}\hspace{0.3em}\rby{
吽=mu}\hspace{0.3em}\rby{
月=mun}\hspace{0.3em}\rby{
楽=musi}\hspace{0.3em}\rby{
多=mute}\hspace{0.3em}\\
\rby{
添=namako}\hspace{0.3em}\rby{
番=nanpa}\hspace{0.3em}\rby{
狂=nasa}\hspace{0.3em}\rby{
道=nasin}\hspace{0.3em}\rby{
山=nena}\hspace{0.3em}\rby{
此=ni}\hspace{0.3em}\rby{
名=nimi}\hspace{0.3em}\rby{
足=noka}\hspace{0.3em}\rby{
お=o}\hspace{0.3em}\\
\rby{
愛=olin}\hspace{0.3em}\rby{
其=ona}\hspace{0.3em}\rby{
開=open}\hspace{0.3em}\rby{
打=pakala}\hspace{0.3em}\\
\rby{
作=pali}\hspace{0.3em}\rby{
支=palisa}\hspace{0.3em}\rby{
米=pan}\hspace{0.3em}\rby{
授=pana}\hspace{0.3em}\rby{
ひ=pi}\hspace{0.3em}\rby{
心=pilin}\hspace{0.3em}\rby{
黒=pimeja}\hspace{0.3em}\rby{
終=pini}\hspace{0.3em}\rby{
虫=pipi}\hspace{0.3em}\rby{
側=poka}\hspace{0.3em}\rby{
箱=poki}\hspace{0.3em}\rby{
良=pona}\hspace{0.3em}\\
\rby{
同=sama}\hspace{0.3em}\rby{
火=seli}\hspace{0.3em}\rby{
皮=selo}\hspace{0.3em}\rby{
何=seme}\hspace{0.3em}\rby{
上=sewi}\hspace{0.3em}\rby{
体=sijelo}\hspace{0.3em}\rby{
丸=sike}\hspace{0.3em}\rby{
新=sin}\hspace{0.3em}\rby{
君=sina}\hspace{0.3em}\rby{
前=sinpin}\hspace{0.3em}\rby{
画=sitelen}\hspace{0.3em}\rby{
知=sona}\hspace{0.3em}\rby{
狸=soweli}\hspace{0.3em}\rby{
大=suli}\hspace{0.3em}\rby{
日=suno}\hspace{0.3em}\rby{
面=supa}\hspace{0.3em}\rby{
甜=suwi}\hspace{0.3em}\\
\rby{
因=tan}\hspace{0.3em}\rby{
許=taso}\hspace{0.3em}\rby{
去=tawa}\hspace{0.3em}\rby{
水=telo}\hspace{0.3em}\rby{
旦=tenpo}\hspace{0.3em}\rby{
語=toki}\hspace{0.3em}\rby{
家=tomo}\hspace{0.3em}\rby{
二=tu}\hspace{0.3em}\\
\rby{
性=unpa}\hspace{0.3em}\rby{
口=uta}\hspace{0.3em}\rby{
鬥=utala}\hspace{0.3em}\\
\rby{
白=walo}\hspace{0.3em}\rby{
一=wan}\hspace{0.3em}\rby{
鴨=waso}\hspace{0.3em}\rby{
力=wawa}\hspace{0.3em}\rby{
无=weka}\hspace{0.3em}\rby{
要=wile}
}

\end{document}
